\section{Design}
\label{sec:design}

\subsection{Overview}
\label{sec:design-overview}

LatchMoE separates replay-stable storage from routing-dependent state changes.
Figure~\ref{fig:architecture} shows the resulting execution model. At
initialization, a CPU-first host store retains expert weights and LatchMoE
allocates fixed device slot banks and persistent logical-to-slot mapping
buffers. During inference, a captured router produces the routed expert IDs.
An eager staging boundary turns those IDs into a placement plan, transfers
missing weights, and updates mapping values in place. A captured expert MLP
then consumes the same slot and mapping allocations on every replay.

\begin{figure*}[t]
  \centering
  \begin{tikzpicture}[
      x=1cm,
      y=1cm,
      font=\small,
      block/.style={draw, rounded corners=1.5pt, align=center,
                    minimum height=9mm, inner sep=4pt},
      stable/.style={block, fill=green!12, draw=green!45!black},
      captured/.style={block, fill=blue!10, draw=blue!55!black},
      dynamic/.style={block, fill=orange!15, draw=orange!70!black},
      context/.style={block, fill=black!4, draw=black!45},
      flow/.style={-{Latex[length=2mm]}, thick},
      control/.style={-{Latex[length=2mm]}, dashed, thick}
    ]
    \node[context, minimum width=22mm] (checkpoint) at (-6.1,2.2) {Model checkpoint};
    \node[stable, minimum width=22mm] (host) at (-3.5,2.2) {CPU-First\\Host Store};
    \node[context, minimum width=22mm] (autoconfig) at (-0.7,2.2) {Capacity\\Configuration};
    \node[stable, minimum width=28mm] (slots) at (2.1,2.2) {Fixed Expert-Slot\\Bank\\stable addresses};
    \node[stable, minimum width=37mm] (mapping) at (5.8,2.2) {Persistent Logical-to-Slot\\Mapping\\stable address, dynamic values};

    \node[captured, minimum width=27mm] (router) at (-5.6,0) {Captured Router};
    \node[dynamic, minimum width=35mm] (staging) at (-1.8,0) {Replay-Boundary Staging\\route, place, transfer, publish};
    \node[captured, minimum width=30mm] (mlp) at (3.4,0) {Captured Expert MLP};
    \node[dynamic, minimum width=35mm] (waves) at (-1.8,-1.5) {Capacity-Bounded\\Wave Executor};
    \node[dynamic, minimum width=32mm] (lifecycle) at (3.4,-1.5) {Versioned Slot Lifecycle\\transfer and compute events};

    \draw[flow] (checkpoint) -- (host);
    \draw[flow] (autoconfig) -- (slots);
    \draw[flow] (router) -- (staging);
    \draw[flow] (host) -- (staging);
    \draw[flow] (staging) -- (mlp);
    \draw[flow] (staging.north east) -- (slots.south west);
    \draw[flow] (slots.south) -- (mlp.north west);
    \draw[flow] (mapping.south) -- (mlp.north east);
    \draw[control] (waves) -- (staging);
    \draw[control] (lifecycle) -- (staging);
    \draw[control] (mlp) -- (lifecycle);

    \node[draw=blue!55!black, dashed, fit=(router)(staging)(mlp)(waves)(lifecycle),
          inner sep=4mm, label={[font=\scriptsize]below:one graph-compatible MoE layer}] {};
  \end{tikzpicture}
  \caption{LatchMoE keeps graph-visible slot and mapping addresses fixed while
  changing their contents and values at an eager replay boundary. Solid arrows
  carry weights or routing data; dashed arrows carry wave and lifecycle
  control. Blue blocks are captured, orange blocks execute eagerly, and green
  blocks are stable allocations.}
  \label{fig:architecture}
\end{figure*}

The design has three core mechanisms. Replay-stable slot virtualization
decouples expert identity from physical allocation. Replay-boundary staging
publishes a new mapping only after transfer readiness. A versioned lifecycle
and a capacity-bounded wave executor make reuse correct when transfer and
compute overlap or when the active set exceeds capacity. The placement-plan
interface orders routed experts but does not embed a prediction policy in the
data plane.

\subsection{Replay-stable expert-slot virtualization}
\label{sec:fixed-slots}

For each managed MoE layer, LatchMoE allocates two contiguous tensors that hold
the layer's fixed set of physical $W_{13}$ and $W_2$ expert slots. The first
dimension indexes physical slots; each slice has the shape and layout expected
by the fused expert kernel. These allocations persist across capture and
replay. A second persistent tensor maps each logical expert ID to a physical
slot ID. Captured computation receives the slot tensors and mapping buffer, so
its graph-visible pointers do not depend on which logical experts are resident.

Each slot records a logical owner, a monotonically increasing version, a state,
and its last-use step. Reassigning a slot removes the old logical owner from
the resident index, installs the new owner, increments the version, and enters
the loading state. A lease is the tuple of slot ID, logical expert, and version.
The version distinguishes two ownership intervals that reuse the same physical
slot and lets asynchronous completions reject stale work.

The mapping is valid only for ready experts. Before the captured MLP consumes
routed IDs, LatchMoE remaps them through the persistent logical-to-slot tensor.
The physical expert count remains the number of allocated slots rather than the
number of logical experts. Thus replacement changes slot contents and mapping
values, while the tensors bound into the captured computation remain fixed.

\subsection{Replay-boundary staging and safe publication}
\label{sec:staging}

The router must run before the runtime knows the active expert set, but the
expert MLP must run after all selected weights are available. LatchMoE exposes
the interval between these operations as an eager staging boundary. The
captured router produces routed expert IDs, the runtime deduplicates that set,
and a placement plan orders the routed experts. The plan must be a permutation
of the actual routed set: it may change staging order, but it cannot add an
unrouted expert or omit a routed expert.

For a cache miss, the transfer engine assigns a slot version, enqueues the
expert's $W_{13}$ and $W_2$ copies on a dedicated H2D stream, and records a
readiness event. Publication follows a fixed order. The runtime first validates
that the slot still matches the lease, installs the readiness event as a
dependency of the consumer stream, changes the slot to ready, and only then
publishes the logical-to-slot mapping. This order prevents captured compute
from observing either incomplete weights or a completion event that belongs to
an older slot owner.

Hits and misses share the same output interface: ready physical slot IDs plus
the persistent mapping tensor. The captured MLP therefore does not encode
whether an expert was already resident, loaded on demand, or prefetched for an
upcoming wave.

\subsection{Versioned transfer--compute lifecycle}
\label{sec:lifecycle}

The slot state machine contains four externally relevant states: empty,
loading, ready, and computing. Empty and ready slots may be selected for a new
owner, while loading and computing slots are not reassignable. Before expert
compute begins, the runtime acquires leases for the active slots and marks
them computing. A completion event protects those leases until the consumer
stream has finished. Release validates that owner and version still match the
recorded lease before returning a slot to the ready state.

This protocol enforces two safety properties. First, transfer completion cannot
make a stale owner ready because readiness checks the version created when the
copy was issued. Second, eviction cannot overwrite an in-flight expert because
a computing slot is excluded from reassignment. Violations raise an error
instead of redirecting execution to eager mode. The fail-closed behavior makes
graph compatibility part of the runtime contract rather than an unobserved
best effort.

\subsection{Capacity-bounded multi-wave execution}
\label{sec:waves}

Decode commonly activates a small expert set, but a prefill request aggregates
routing decisions over many tokens. The union of routed experts can exceed the
slot capacity even though each token selects only a few experts. LatchMoE
handles this overflow by partitioning routed token--expert pairs into waves.
Each wave contains no more distinct missing experts than its staging bank can
hold. Hit-only work may be scheduled separately so that a ready main-slot
expert need not be copied into a temporary wave buffer.

Wave execution preserves the router's assignments. LatchMoE retains every
routed pair's original flattened top-$k$ position, computes wave-local expert
outputs, concatenates those outputs, reconstructs one global native unpermute
index, and invokes the ordinary token-combine operation once with the original
top-$k$ weights and output shape. It rejects incomplete or duplicate pair
coverage. The single final combine avoids the numerical reordering introduced
by independently combining and accumulating each wave.

Multiple stage banks form a software pipeline. While the current wave computes,
the transfer stream may populate another bank for the next wave. The scheduler
can order waves by their required transfer bytes, but it cannot change pair
membership or final combination semantics. A stage bank is reused only after
its compute lease completes, so the same versioned lifecycle protects both the
persistent decode slots and transient prefill slots.

\subsection{Policy-neutral placement}
\label{sec:policy-neutral}

LatchMoE defines the data-plane constraints for a placement plan: target layer,
complete routed set, and a valid ordering. The default behavior preserves the
router's first-seen order. An external controller may instead prioritize
predicted reuse or transfer cost, but it must return a permutation of the
routed experts and obey the same capacity and lifecycle checks. Separating
this interface from slot management allows placement policies to share one
graph-compatible execution substrate and keeps policy quality outside the
correctness argument of this paper.
